%!TEX root = ../main.tex
\section{Methods}
\label{sec:methods}

The five final submissions came from three model families. Each was designed
to capture a different view of short-term cold-related acoustics while keeping
the final classifier small.

\begin{table*}[t]
\centering
\caption{The three submitted model families. All final decision models are
low-capacity even when their input representation is learned.}
\label{tab:families}
\small
\begin{tabular}{p{2.5cm}p{6.4cm}p{4.0cm}c}
\toprule
Family & Acoustic view & Classifier/fusion & Sub. \\
\midrule
Compact handcrafted & Seven gain-invariant energy/pause features (G4) and 168 CQT summaries (G9) & Balanced logistic regression; fixed equal standardized-logit fusion & 1--2 \\
Trajectory & Depth-2 log-signatures of MFCC+time paths over the full clip and two halves & Balanced logistic regression; fixed equal fusion with G4 & 3 \\
Global/local deep & HuBERT layer 3 + eGeMAPS global statistics and a three-channel acoustic-image CNN & PCA+Ridge and lightweight CNN; fixed equal fusion, optionally diagonal CORAL & 4--5 \\
\bottomrule
\end{tabular}
\end{table*}

\subsection{Compact handcrafted fusion (Submissions 1--2)}

The first branch, \textbf{G4}, represents energy and pausing with seven
gain-invariant descriptors. Frame-level root-mean-square energy and
voiced/unvoiced/silence regions from pYIN yield a low-energy ratio, energy
slope, three relative dB contrasts, long-pause rate and median silence
duration. We excluded four absolute-loudness variables to reduce dependence on
recording gain.

The second branch, \textbf{G9}, summarizes a constant-Q transform (CQT). We
use 84 bins at 12 bins per octave from approximately C1, covering seven
octaves. The temporal mean and standard deviation of each log-amplitude bin
produce 168 features. Unlike a fixed-width mel front end, the logarithmic
resolution preserves detail around fundamental frequency and low-frequency
resonance; earlier URTIC work also found CQT representations useful
\cite{cai2017cold}.

Each branch uses a StandardScaler and class-balanced logistic regression.
Let $\ell_4$ and $\ell_9$ be their logits and let $z(\ell)$ denote
standardization with mean and variance fitted on the training pool. The final
score is deliberately parameter-free,
\begin{equation}
 s_{4,9}(x)=\tfrac{1}{2}\big[z(\ell_4(x))+z(\ell_9(x))\big].
 \label{eq:fusion}
\end{equation}
Submission 1 used the fixed boundary $s_{4,9}\geq0$. Submission 2 used a
threshold calibrated from fitting-side OOF predictions. We avoided a learned
fusion weight because its sign and magnitude were unstable across folds.

\subsection{Path log-signatures (Submission 3)}

Global summary statistics discard the order in which the voice changes. For
each clip we therefore form a path from 20 frame-level MFCCs sampled every
10~ms plus normalized time. A depth-2 log-signature contains first-order total
increments and second-order ordered interactions between channel pairs. For a
21-dimensional path this gives 231 coefficients. We compute one signature on
the full interval and one on each dyadic half, producing $3\times231=693$
features. Splitting into halves prevents a short cough or throat-clear near one
end from being averaged across the entire clip. A class-balanced logistic
regression operates on the signatures; its standardized score is averaged
equally with G4.

\begin{figure*}[t]
\centering
\resizebox{\textwidth}{!}{%
\begin{tikzpicture}[
  box/.style={draw,rounded corners,align=center,minimum height=12mm,
              text width=27mm,font=\small},
  wide/.style={box,text width=36mm},
  score/.style={box,text width=33mm},
  arr/.style={-{Latex[length=2mm]},thick}]
\node[box] (wav) {waveform\\16-kHz mono};

\node[wide,right=8mm of wav,yshift=12mm] (manner) {pYIN + RMS manner\\labels at 50~Hz};
\node[wide,right=8mm of manner] (g4) {G4: 11 scalars, 7 retained\\(gain-invariant subset)};
\node[box,right=8mm of g4] (lr4) {scaler +\\balanced LR};
\node[box,right=8mm of lr4] (z4) {$z$-score\\(fit-side statistics)};

\node[wide,right=8mm of wav,yshift=-12mm] (mfcc) {MFCC path, 20 channels\\at 100~Hz, time appended};
\node[wide,right=8mm of mfcc] (gsig) {$G_{\mathrm{sig}}$: depth-2 log-signature\\whole path + two halves (693-d)};
\node[box,right=8mm of gsig] (lrs) {scaler +\\balanced LR};
\node[box,right=8mm of lrs] (zs) {$z$-score\\(fit-side statistics)};

\node[score,right=10mm of z4,yshift=-12mm] (fuse) {fused score\\$\tfrac12(z_{G4}+z_{G\mathrm{sig}})$\\predict C if score $\geq0$,\\else NC};

\draw[arr] (wav.east) -- ++(4mm,0) |- (manner.west);
\draw[arr] (wav.east) -- ++(4mm,0) |- (mfcc.west);
\draw[arr] (manner)--(g4); \draw[arr] (g4)--(lr4); \draw[arr] (lr4)--(z4);
\draw[arr] (mfcc)--(gsig); \draw[arr] (gsig)--(lrs); \draw[arr] (lrs)--(zs);
\draw[arr] (z4.east) -- ++(4mm,0) |- (fuse.west);
\draw[arr] (zs.east) -- ++(4mm,0) |- (fuse.west);
\end{tikzpicture}%
}
\caption{Submission 3: the gain-invariant G4 branch and the MFCC
log-signature branch are trained independently, standardized using fitting-side
statistics and fused with fixed equal weights. This was the best hidden-Test
system (UAR 0.6582).}
\Description{Two-branch audio classification pipeline. Gain-invariant
energy-and-pause descriptors and MFCC path log-signatures feed separate
balanced logistic regressions whose standardized scores are averaged.}
\label{fig:signature_pipeline}
\end{figure*}

\subsection{Global/local multi-view system (Submissions 4--5)}

The \textbf{global view} concatenates 4,608-dimensional pooled HuBERT Base
layer-3 statistics with 88 eGeMAPSv02 functionals~\cite{eyben2016egemaps},
giving 4,696 features. A variance filter, StandardScaler and PCA (256
components), all fitted on the source side, precede a class-balanced Ridge
classifier ($\alpha=10$). Layer 3 was the strongest explored HuBERT layer on
the internal acoustic-sensitivity screen.

The \textbf{local view} treats each clip as a three-channel acoustic image:
log-mel spectrum, MFCCs and $\Delta$MFCCs, each represented as $40\times512$
and stacked into a $3\times40\times512$ tensor. After per-channel
standardization, a lightweight CNN with BatchNorm, GELU, dropout 0.35, global
average pooling and one fully connected output learns local time--frequency
textures that global pooling removes.

The Ridge score $S_1$ and CNN logit $S_2$ are independently standardized using
fitting-side statistics and averaged with fixed 50/50 weights. Submission 4
uses the historical threshold $t^*=0.3913$, selected on the transferred
$k=210$ pseudo-speaker validation split. This records the submitted procedure;
Sec.~\ref{sec:data_eval} explains why that within-Development split is not a
defensible speaker-disjoint evaluation. Submission 5 additionally applies
diagonal CORAL~\cite{sun2016coral}: feature means and marginal variances are
aligned between the labeled fitting distribution and unlabeled Test
recordings. No Test labels are used, but the procedure is transductive and is
reported as such.

\begin{figure*}[t]
\centering
\resizebox{0.94\textwidth}{!}{%
\begin{tikzpicture}[
  input/.style={draw=blue!70!black,fill=blue!6,rounded corners,
                align=center,text width=66mm,minimum height=10mm,font=\small},
  global/.style={draw=green!45!black,fill=green!5,rounded corners,
                 align=center,text width=68mm,minimum height=30mm,font=\small},
  local/.style={draw=violet!75!black,fill=violet!5,rounded corners,
                align=center,text width=68mm,minimum height=30mm,font=\small},
  fusion/.style={draw=orange!80!black,fill=orange!6,rounded corners,
                 align=center,text width=96mm,minimum height=12mm,font=\small},
  arr/.style={-{Latex[length=2mm]},thick}]
\node[input] (audio) {\textbf{Audio input: 16-kHz mono waveform}};

\node[global,below left=8mm and 5mm of audio] (g) {
\textbf{Stream 1: global acoustic summary}\\[1mm]
HuBERT Base L3 (4,608-d) + eGeMAPS (88-d)\\
$\downarrow$ concatenate (4,696-d)\\
$\downarrow$ variance filter $\rightarrow$ scaler $\rightarrow$ PCA (256)\\
$\downarrow$ balanced Ridge ($\alpha=10$) $\rightarrow S_1$};

\node[local,below right=8mm and 5mm of audio] (l) {
\textbf{Stream 2: local time--frequency dynamics}\\[1mm]
Log-mel + MFCC + $\Delta$MFCC (each $40\times512$)\\
$\downarrow$ stack ($3\times40\times512$), per-channel standardize\\
$\downarrow$ Conv2D $\rightarrow$ BN $\rightarrow$ GELU $\rightarrow$ dropout 0.35\\
$\downarrow$ global average pool $\rightarrow$ FC $\rightarrow S_2$};

\node[fusion,below=9mm of $(g.south)!0.5!(l.south)$] (f) {
$Z_i=(S_i-\mu_i)/\sigma_i$ using fitting-side statistics
$\quad\rightarrow\quad Z_{\mathrm{fused}}=0.5Z_1+0.5Z_2$\\
predict Cold if $Z_{\mathrm{fused}}\geq t^*=0.3913$, otherwise Non-Cold};

\draw[arr] (audio.south) -- ++(0,-3mm) -| (g.north);
\draw[arr] (audio.south) -- ++(0,-3mm) -| (l.north);
\draw[arr] (g.south) |- (f.west);
\draw[arr] (l.south) |- (f.east);
\end{tikzpicture}%
}
\caption{Submissions 4--5: Akshat's global/local multi-view system. The two
views are fused with fixed equal weights; Submission 5 additionally applies
diagonal CORAL. It achieved the best hidden-Test accuracy (0.7708) and F1
(0.5800), while its UAR (0.6567) was close to the signature system. The shown
threshold is the historical submitted policy, not a corrected evaluation
threshold.}
\Description{Two-stream audio classification pipeline. Pooled HuBERT and
eGeMAPS features feed a Ridge classifier, while stacked log-mel, MFCC and
delta-MFCC features feed a convolutional network. Their standardized scores
are averaged and thresholded.}
\label{fig:multiview_pipeline}
\end{figure*}

\subsection{Exploratory models and controls}

Before the final submissions, we tested frozen WavLM/HuBERT pooled-statistic
heads, audit-derived layer weighting, learned and constrained late fusion,
waveform splicing, embedding mixup, speaker-masked supervised contrastive
learning, gradient-reversal speaker adversaries, calibration variants and
test-time augmentation. These experiments were valuable mainly as negative
controls. Because most used the historical within-Development boundary, their
absolute UARs are exploratory; only comparisons made on a common boundary and
the failure mechanisms remain interpretable.
